\section{Conclusiones y Trabajos Futuros}

\subsection{7.1 Conclusiones}

Particularmente llamativo es el hecho de que avances aparentemente modestos en planificación y control puedan generar mejoras sustanciales cuando se aplican al ensamblaje robótico in-orbit. Este trabajo ha presentado un framework completo que integra arquitectura modular multi-brazo, planificación híbrida optimizada, control de posición-fuerza y un gemelo digital bidireccional. Los resultados obtenidos demuestran reducciones significativas en tiempo de ensamblaje (hasta 50\% en casos estudiados), consumo energético y riesgo operativo global.

Revisiones recientes del campo refuerzan la relevancia de estos desarrollos al identificar precisamente brechas en eficiencia y autonomía que el enfoque propuesto ayuda a cerrar \cite{Wang2025_OASR}. Más allá de los números, el sistema representa un paso concreto hacia infraestructuras espaciales que no solo se lanzan, sino que se construyen y reparan en órbita. Esto resulta especialmente valioso para exploración profunda, donde cada kilogramo y cada minuto cuentan de forma decisiva.

Los hallazgos confirman que la combinación de coordinación multi-robot con predicción en tiempo real mediante gemelos digitales ofrece un camino viable para superar limitaciones históricas del estado del arte. No obstante, cabe matizar que estos avances se basan fundamentalmente en simulaciones de alta fidelidad. La verdadera validación vendrá con pruebas en condiciones orbitales reales.

\subsection{7.2 Recomendaciones y Trabajos Futuros}

Uno de los desafíos persistentes que surge al proyectar estos sistemas hacia el futuro radica en escalar la autonomía real bajo condiciones impredecibles y recursos limitados. El estado actual de capacidades ISAM sugiere que todavía queda un trecho importante para alcanzar operaciones sostenidas y completamente autónomas en el espacio profundo \cite{NASA_ISAM2025}.

Entre las líneas prioritarias se encuentra el avance hacia configuraciones multi-robot distribuidas con mayor número de agentes cooperativos. Esto exige nuevos protocolos de coordinación descentralizada y mecanismos robustos de reallocación de tareas ante fallos. Igualmente prometedora resulta la integración más profunda de técnicas de machine learning para control adaptativo y estimación de estado en entornos no estacionarios \cite{Alizadeh2024_Survey}.

Otras direcciones relevantes incluyen la incorporación de manufactura aditiva in-situ durante el ensamblaje, el desarrollo de materiales inteligentes con sensores embebidos y la optimización multi-objetivo que considere explícitamente la longevidad de la misión completa. También sería valioso explorar arquitecturas híbridas humano-robot para etapas de transición en misiones tripuladas a la Luna o Marte.

En definitiva, el trabajo aquí desarrollado establece bases sólidas, pero representa solo un escalón en un camino mucho más amplio. La convergencia entre robótica avanzada, gemelos digitales y aprendizaje automático abre perspectivas reales para construir en el espacio de manera similar a como construimos en la Tierra: de forma incremental, eficiente y adaptable. El siguiente paso lógico pasa por hardware en órbita. Ese será el verdadero banco de pruebas.